\chapter{Conclusion}
\label{ch:conclusion}

The thesis asked whether choosing demonstrations can help a relatively small model, at 7B, follow them over its own priors under controlled distribution shift. To answer it, we built a synthetic tabular testbed in which covariate shift, spurious reversal and concept shift can each be switched on alone, set the model's semantic prior against the rule by naming the features, and tested every confirmatory claim against a contrast family frozen before its data were seen.

For Qwen2.5-7B-Instruct with 8 numeric demonstrations, none of the selection methods we tried did so. The feature names decided most of what the model predicted, more than any shift of the data. The hand-designed strategies differed little from balanced random demonstrations; the one clear effect is that neutralising the shortcut fares better under its reversal than under covariate shift. A learned selector, the only method with no weak environment on a stand-in learner, did no better than balanced random demonstrations on the model, though in an exploratory study it did help the base model, whose priors pull less. Inside the model, the answers form in the last few layers, and the names decide them there. The prior is present from the start and the demonstrated rule hardly at all, and a single steering direction moves the answers without making them more correct.

Selection chooses evidence; in our experiments, it did not make the instruction-tuned model read evidence it largely ignored. Within the scale we could study, this appears to be the binding constraint. The testbed we have constructed here, including the pre-registered families and the diagnosis, gives future work, at a larger scale and with the model's own training signal, a controlled place to find out when it stops being a constraint.
